\begingroup
\makeatletter\tagsleft@true\let\veqno\@@leqno\makeatother
\ArithmeticHead{268}{L}{ANALYSIS. NACHLASS.}
\SourceLeaf{275}
\pdfbookmark[1]{Erläuterungen zu [IV.]}{x1-agm-erlaeuterung-iv}
\begin{center}\textls[120]{Erläuterungen zu [IV.], S. 184—193.}\end{center}
\par\indent Die ersten 8 Seiten der Scheda Ac (begonnen November 1799) enthalten Aufzeichnungen über die analytische Geometrie des Raumes und die sphärische Trigonometrie. Nur auf S. 7 finden sich zwischen diesen geometrischen Notizen, aber in sorgfältiger Schrift und von Linien eingerahmt, drei Zeilen, mit den Reihenentwicklungen [1], [2], [3] (oben S. 184) für $\frac\varpi\pi$, von denen die erste und dritte schon Werke III, S. 423, art. [15.] abgedruckt sind. Die erste Entwicklung [1] stimmt mit der auf der Titelrückseite des \textsc{Leiste} stehenden und oben S. 169 wiedergegebenen Reihe überein. Ebenda ist auch die Reihe [2] angedeutet und mit ihrer Hilfe der Wert von $\frac\varpi\pi$ auf 15 Dezimalstellen berechnet. Dieselbe Reihe [2] ist auch im \textit{Tagebuch}, Nr. 91a mit der Datierung Juli 1798 aufgezeichnet. Mit der S. 8 schließen die geometrischen Betrachtungen ab und es folgen auf S. 9 zunächst die beiden Reihen [4], [5] für $\frac\varpi\pi-\frac\pi\varpi$, von denen die zweite Werke III, S. 423, art. [15] abgedruckt ist. Dann folgt die Aufzeichnung art. [2.], zu der man den Schluß des art. [2.] des vorigen Abschnitts [III.], S. 183, ferner den Brief an \textsc{Bessel} vom 3. September 1805 (siehe oben Briefwechsel, [3.], S. 238) vergleichen kann. Zu diesen Entwicklungen dürfte \textsc{Gauss} durch Untersuchungen, die der theoretischen Astronomie angehören, veranlaßt worden sein. In [6] bedeutet $c$ den Kosinus eines Winkels, $c2,c3,\ldots$ sind die Kosinus seiner Vielfachen. Bemerkenswert ist die elegante Kettenbruchentwicklung für den Quotienten der beiden ersten Koeffizienten $\frac{A'}A$. Dieser Quotient, der auch in der Aufzeichnung [25]—[28], S. 188, betrachtet wird, steht wie z. B. die Entwicklungen [28] zeigen, in engem Zusammenhang mit der im \textsc{Leiste} vorkommenden Funktion $\mathbf M$ (vergl. insbesondere Abschnitt [II.] Gl. [13] und [28])\textsuperscript{*)}. Der Zusammenhang dieser Entwicklungskoeffizienten mit dem agM. wird in wenig abweichender Form in den Gln. [23], [24] S. 188 hervorgehoben.
\par\indent Während die bisher betrachteten Notizen über das agM. vorwiegend formaler Natur waren, führen die Aufzeichnungen der artt. [3.] und [4.] in die Tiefen dieser Theorie ein. Auf den Weg, der \textsc{Gauss} zu der wichtigen Quotientendarstellung [7] geführt haben mag, wirft die Gleichung [10], S. 186, ein helles Licht, wenn man sie mit den Leisteformeln des art. [1.] von Abschnitt [II.] S. 177 vergleicht.
\par\indent Das Verhalten der Funktion $M(1,x)$, die \textsc{Gauss} mit $Mx$ bezeichnet, für große Werte von $x$ wird in erster Annäherung durch die Gleichung
\begin{equation*}\tag*{$(\alpha)$}\lim_{x\to\infty}\frac1xM(1,x)\log4x=\frac\pi2\end{equation*}
\noindent gegeben, die mit der Gleichung (14a) unseres \textit{Abrisses} (S. 255) übereinstimmt. Diese Gleichung kommt zwar in den hier zu besprechenden Aufzeichnungen nicht ausdrücklich vor, sie liegt aber, wie wir sehen werden, der Quotientendarstellung [7] zugrunde und spielt auch in späteren Entwicklungen von \textsc{Gauss} zum agM. namentlich bei der Ableitung des »Schönen Lehrsatzes«\textsuperscript{**)} eine wichtige Rolle, wo auch ein Weg für ihre Ableitung angedeutet ist; wir haben im Artikel 2. des \textit{Abrisses} diese Andeutungen ausgeführt. Es kommt dabei die Beziehung des agM. zum Ellipsenquadranten zur Geltung, indem nämlich durch diese Beziehung die Anwendung der Methoden ermöglicht wird, die \textsc{Euler} in der oben (S. 254, Fußnote) genannten Abhandlung \textit{Animadversiones etc.} der Opuscula varii argumenti II angegeben hat. In dieser Abhandlung untersucht \textsc{Euler} das Verhalten des Ellipsenquadranten $q$, wenn die große Achse den festen Wert 1
\Needspace{4\baselineskip}
\par\smallskip\noindent\rule{.25\linewidth}{.4pt}\par
{\small
\par\indent *) Man vergl. auch Werke III, S. 371, und \textsc{Gauss}' Brief an \textsc{Schumacher} vom April 1816, siehe oben Briefwechsel [7.], S. 247, ferner die allgemeinen Formeln, Werke III, S. 128, 129, (\textit{Disquisitiones circa seriem etc.} 1812).
\par\indent **) Siehe Abschnitt [IX.] art. [4.], S. 220.
\par}
\clearpage
\fancyhead{}\fancyhead[R]{\small269}\fancyhead[C]{\scriptsize BEMERKUNGEN ZUM AGM. ABSCHNITT IV.}
\SourceLeaf{276}
\noindent behält, während die kleine Achse $p$ sich der Null annähert. Zunächst folgt das auch geometrisch evidente Resultat, daß
\[\lim_{p\to0}q=1\]
\noindent sein muß. Die Betrachtung der linearen Differentialgleichung zweiter Ordnung, der $q$ als Funktion von $p$ genügt (a. a. O. S. 147):
\[p(1-p^2)\frac{d^2q}{dp^2}-(1+p^2)\frac{dq}{dp}+pq=0,\]
\noindent leitet \textsc{Euler} zu dem Ansatz (a. a. O. S. 155)
\[q=1+Ap^2+Bp^4+\,.\,.-\left(\frac12p^2+\beta p^4+\,.\,.\right)\log p.\]
\noindent Für $A$ findet \textsc{Euler} zunächst induktiv durch numerische Rechnung den Wert
\[A=\log2-\frac14,\]
\noindent den er dann durch Vergleichung des Ellipsenbogens mit einem Parabelbogen bestätigt. In der von \textsc{Borchardt} ausgearbeiteten und im Bande I der Werke \textsc{Jacobis} auszugsweise abgedruckten Vorlesung von \textsc{Jacobi}, wird (S. 522—525) die \textsc{Euler}sche Methode auf das Integral erster Gattung angewandt. In der ursprünglichen \textsc{Borchardt}schen Ausarbeitung, die von dem abgedruckten Texte nicht unbeträchtlich abweicht, lautet die betreffende Stelle, wie folgt:
\par\indent »Man erhält als Endwert für $\varkappa=0$
\[K'=\log\frac4{\varkappa}\ \text{\textsuperscript{*})]}.\]
\noindent Dieses Resultat hat zuerst \textsc{Euler} gefunden und veröffentlicht in dem Werke Opuscula varii argumenti, welches man häufig auf Auktionen bekommt; es findet sich auch im \textsc{Legendre} [\textit{Exercices de calcul intégral} I, 1811, § 72ff.]. Die Schwierigkeit seiner Herleitung bestand nicht in der Form des Resultats, da man schon lange wußte, daß man bei solchen Entwicklungen, welche Potenzen enthalten, auch einen logarithmischen Term beifügen muß; ja man konnte in unserem Falle sogar leicht finden, daß der Endwert $=\log\frac n{\varkappa}$ war, aber der numerische Wert von $n$ war sehr schwer zu ermitteln, was man bei \textsc{Legendre} nachsehen kann, der zwar den Wert $n=4$ findet, aber keinen stringenten Beweis dafür gibt.«
\par\indent Wie man sofort erkennt, ist die Gleichung \textsc{Jacobis} mit der oben angegebenen Gleichung $(\alpha)$ gleichwertig. Da \textsc{Gauss} die Beziehung zwischen dem agM. und dem Ellipsenquadranten gekannt hat (siehe Abschnitt [II.] Gln. [14], [15], S. 178), so könnte er die Gleichung $(\alpha)$ bezw. die Gleichung
\begin{equation*}\tag*{$(\beta)$}\lim_{p\to0}M(1,p)\log\frac4p=\frac\pi2\end{equation*}
\noindent in folgender Weise unmittelbar aus dem \textsc{Euler}schen Resultat abgeleitet haben. Nehmen wir $a=1$, so folgt aus den Formeln [14] (siehe die berichtigte Formel $[14]'$ S. 265) und [35] des Abschnitts [II.] S. 180, wenn man $\nu^2=1-p^2$ setzt, die Gleichung
\[\frac1{M(1,p)}=K(\nu)=\frac2\pi\left(\frac{1-p^2}{p}\frac{dq}{dp}+q\right).\]
\Needspace{5\baselineskip}
\par\smallskip\noindent\rule{.25\linewidth}{.4pt}\par
{\small\par\indent *) Hier ist $K'=\displaystyle\int_0^{\frac\pi2}\frac{d\varphi}{\sqrt{1-(1-\varkappa^2)\sin^2\varphi}}$.\par}
\clearpage
\ArithmeticHead{270}{L}{ANALYSIS. NACHLASS.}
\SourceLeaf{277}
\par\indent Nun ist nach \textsc{Euler}
\[q=1+\left(\log2-\frac14\right)p^2-\frac12p^2\log p+p^4[\ldots],\]
\noindent also
\[\frac{dq}{dp}=(2\log2-1)p-p\log p+p^3[\ldots]\]
\noindent und
\[\frac{1-p^2}{p}\frac{dq}{dp}+q=2\log2-\log p+p^2[\ldots],\]
\noindent woraus die Gleichung $(\beta)$ ohne weiteres hervorgeht. Daß \textsc{Gauss} ursprünglich diese Herleitung angewandt hat, ist darum sehr wahrscheinlich, weil sie sich an die im \textsc{Leiste} aufgezeichneten Formeln des Abschnitts [II.] ganz \textls[120]{unmittelbar} anknüpfen läßt.
\par\indent Dagegen kann der von \textsc{Gauss} in der späteren Aufzeichnung des Abschnittes [IX.] art. [4.] angedeutete Weg mehr als eine \textls[120]{Übertragung} des von \textsc{Euler} für den Ellipsenquadranten vorgezeichneten Verfahrens auf das agM. bezw. auf die Funktion $K(\nu)$ bezeichnet werden; auch dafür standen \textsc{Gauss} schon hier alle erforderlichen Mittel zur Verfügung, indem er aus den Differentialrelationen [37] des Abschnitts [II.] nur die lineare Differentialgleichung zweiter Ordnung für $K(\nu)$ herzuleiten brauchte, um Schritt für Schritt dem Vorgange \textsc{Eulers} folgen zu können.
\par\indent Jedenfalls können wir sagen, \textls[120]{daß} \textsc{Gauss} \textls[120]{das Verhalten von} $M(1,x)$ \textls[120]{für große Werte von} $x$ \textls[120]{mit Hilfe des Ellipsenquadranten erforscht hat und zwar auf einem Wege, der ihm durch die} \textsc{Euler}\textls[120]{sche Abhandlung »}\textit{Animadversiones etc.}\textls[120]{« vorgezeichnet war.}
\begin{center}*\hspace{2cm}*\\[-.7ex]*\end{center}
\par\indent Der weitere Fortgang von \textsc{Gauss}' Untersuchungen liegt jetzt klar zu Tage.
\par\indent Es ist nach der Formel [3] im \textsc{Leiste} (Abschnitt [II.], S. 177)
\begin{equation*}\tag*{$[3]'$}B=4z-16z^3+56z^5-\ldots;\end{equation*}
\noindent setzt man hierin
\[B=\frac1x,\quad z=\frac1{\mathfrak z},\]
\noindent so kommt die Formel
\begin{equation*}\tag*{$[10]'$}x=\frac1{\dfrac4{\mathfrak z}-\dfrac{16}{\mathfrak z^3}+\dfrac{56}{\mathfrak z^5}-\dfrac{160}{\mathfrak z^7}-\ldots}\end{equation*}
\noindent zum Vorschein, die mit [10], S. 186 übereinstimmt; wir haben nur der Deutlichkeit wegen $\mathfrak z$ an Stelle des von \textsc{Gauss} benutzten $z$ geschrieben, weil später in [22], S. 188 und von [32], S. 190 an) $z$ in derselben Bedeutung wie im \textsc{Leiste} benutzt wird.
\par\indent Aus $[10]'$ ergibt sich
\[\log\mathfrak z=\log\left(4x-\frac1x-\frac9{32}\frac1{x^3}-\ldots\right),\]
\noindent so daß die Gleichung $(\alpha)$ auch in der Form
\[\lim_{x\to\infty}\frac1xM(1,x)\log\mathfrak z=C\]
\noindent geschrieben werden kann. Dies veranlaßt nun zu setzen
\[M(1,x)=\frac{C(x-\alpha x^{-1}-\beta x^{-3}-\ldots)}{\log\mathfrak z};\]
\noindent die Koeffizienten $\alpha,\beta,\ldots$ lassen sich numerisch bestimmen, ihre Werte stehen im art. [3.], S. 186.
\clearpage
\fancyhead{}\fancyhead[R]{\small271}\fancyhead[C]{\scriptsize BEMERKUNGEN ZUM AGM. ABSCHNITT IV.}
\SourceLeaf{278}
\par\indent Jetzt macht \textsc{Gauss} die »Observatio« [9]. Diese leitet ihn zu der Vermutung, daß der Zähler »omisso factore $C$« auch allgemein ein agM. sein könnte, und da der reziproke Wert des agM. einfacher zu behandeln ist, als das agM. selbst, so setzt er in [13] den Zähler $=\frac CQ$ und findet für $Q$ die Entwicklung [14], aus der er gemäß Abschnitt [III.] art. [2.], S. 182 die Aussage [15] erschließt. Nun macht er ähnlich wie schon für $x=\sqrt2$ im \textsc{Leiste} (siehe oben S. 149, art. [8.]) den Übergang zur Integraldarstellung, schreibt aber zunächst, wohl infolge eines Rechenfehlers, das unrichtige
\[\int_0^1\sqrt{\frac{1-r^2}{x^2-r^2}}\,dr,\]
\noindent neben das er erst später (mit ersichtlich anderen Schriftzügen und anderer Tinte) berichtigend das Integral erster Gattung setzt\textsuperscript{*)}. Die Reihenentwicklung [14], die genau mit der Funktion $\varphi x$ der Aufzeichnung [III.] art. [1.] übereinstimmt, läßt jetzt $Q$ als agM. erkennen, was in der Gleichung [20] ausgesprochen ist. »Numerator pro $\frac1k$« bedeutet dort soviel, wie der Zähler des Ausdrucks in [7] für $x=\frac1k$. Das »demonstr[atum]« deutet darauf hin, daß es sich um den Beweis einer früher gehegten Vermutung handelte (vergl. oben), die unerklärliche Buchstabenverbindung GALEN (verschnörkelte Großbuchstaben) erinnert an das »vicimus GEGAN« der Tagebuchnotiz Nr. 43 vom 21. Oktober 1796.
\par\indent Die jetzt gewonnene \textls[120]{allgemeine} Einsicht, daß der reziproke Wert des agM. das vollständige Integral erster Gattung ist, liefert nicht nur aufs neue den Beweis der in der Tagebuchnotiz vom 30. Mai 1799 gemachten Bemerkung (siehe oben S. 260), sondern muß als ausschlaggebend dafür angesehen werden, daß \textsc{Gauss} von nun ab seine Aufmerksamkeit dem allgemeinen elliptischen Integral \textls[120]{erster} Gattung als der wahren Verallgemeinerung des lemniskatischen Integrals zuwendet. Die »Peripheria Ellipseos« tritt in der Gleichung [22] noch einmal auf. Die kleine Zahlenrechnung zeigt, daß \textsc{Gauss} hier die Formel [4] des \textsc{Leiste} (Abschnitt [II.], S. 177) nachgerechnet hat, indem er in die Entwicklung [1] (S. 177) den Wert [3] einsetzt. Er begeht aber wieder einen Rechenfehler, denn statt $-176$ muß richtig $+551$ stehen, wodurch statt $-64z^3$ der richtige Wert $+663z^3$ zum Vorschein kommt.
\par\indent Mit der Formel [22] ändert sich die Bedeutung des Buchstaben $z$; von nun ab ist $z$ der reziproke Wert der in [10] mit demselben Buchstaben bezeichneten Größe (unseres $\mathfrak z$), hat also dieselbe Bedeutung wie in den Formeln des \textsc{Leiste}, Abschnitt [II.]. Nach einigen Zwischenrechnungen drückt \textsc{Gauss} in [32] die Größe $z$ durch $s=\frac1x$ aus, die logarithmische Derivierte dieses Ausdrucks wird schon in [29] betrachtet. Dabei hat $s$ dieselbe Bedeutung wie das $B$ in der Leisteformel [3] Abschnitt [II.], S. 177, so daß also [32] durch Inversion der berichtigten Reihe [3] bezw. der in $[10]'$, S. 270 auftretenden Reihe, für $x=\frac1s$, $\mathfrak z=\frac1z$ hervorgeht. Nun bildet er die Funktion $P(s)$, die nach der Reihenentwicklung [36] (vergl. [33] im \textsc{Leiste}, Abschnitt [II.], S. 180) und nach [37] (vergl. [35] im \textsc{Leiste}, Abschnitt [II.], S. 180) durch die Gleichung
\begin{equation*}\tag*{$[37]'$}M(1,\sqrt{1-s^2})=\frac1{P(s)}\end{equation*}
\noindent definiert werden kann. Der in [37] benutzte Buchstabe $c$ hat nämlich die Bedeutung $\sqrt{1-s^2}$, wobei man etwa an cosinus oder an die \textsc{Gauss} geläufige Bezeichnung $c=\sqrt{a^2-b^2}$ für $a=1$, $b=s$ zu denken hat. Setzt man in die Reihe [36] für $s$ seinen Ausdruck aus $[10]'$
\begin{equation*}\tag*{$[10]''$}s=4z-16z^3+56z^5-160z^7-\ldots\end{equation*}
\Needspace{4\baselineskip}
\par\smallskip\noindent\rule{.25\linewidth}{.4pt}\par
{\small\par\indent *) Der Abdruck gibt an dieser Stelle (S. 187) die Anordnung der Handschrift genau wieder.\par}
\clearpage
\ArithmeticHead{272}{L}{ANALYSIS. NACHLASS.}
\SourceLeaf{279}
\noindent ein, so ergibt sich für $P(s)$ die Darstellung [33], deren Umformung in die Summe von zwei Quadraten [34] sehr nahe liegt. Die Gleichung [35] dürfte \textsc{Gauss} dann mit Hilfe der ihm bekannten Formel (\textsc{Leiste}, [6] Abschnitt [II.], S. 177)
\begin{equation*}\tag*{$[6]'$}sP(s)=s+\frac14s^3+\frac14\,.\,\frac9{16}s^5+\ldots=(2z^{\frac12}+2z^{\frac92}+\ldots)^2\end{equation*}
\noindent abgeleitet haben. Die Formel [38] und die ihr äquivalente [39] (in dieser hat $c$ wieder die Bedeutung $\sqrt{1-s^2}$) ist nun nichts anderes als die Quotientendarstellung aus [7] in den neuen Bezeichnungen; wir können sie in der Form
\begin{equation*}\tag*{$[38]'$}\log z=-\frac\pi2\frac{M(1,\sqrt{1-s^2})}{M(1,s)}\end{equation*}
\noindent schreiben, in der sie dann unmittelbar mit der Hauptformel (38) unseres \textit{Abrisses} (S. 256) übereinstimmt. Wir werden ihr auf S. 34 der Scheda Ac (Abschnitt [V.] art. [5.], S. 197) wieder begegnen.
\par\indent Wir sehen, daß sich \textsc{Gauss} durch diese Untersuchungen in den vollen Besitz der Theorie des agM. gesetzt hat. Während er im \textsc{Leiste} die Größe $z$ nur in formaler Weise, nämlich durch den Reihenansatz [6], S. 177 definierte, erzielt er hier die genuine analytische Definition $[38]'$ von $z$ durch das agM., d. h. er hat jetzt das geleistet, was im Artikel 3. unseres \textit{Abrisses} dargestellt ist. Endlich eröffnet ihm der Zusammenhang zwischen dem agM. und dem vollständigen elliptischen Integral erster Gattung den Zugang zu der allgemeinen Theorie der elliptischen Funktionen, indem die Formel $[38]'$ die Größe $-\frac2\pi\log z$ direkt durch den Periodenquotienten darstellt\textsuperscript{*)}.
\par\indent Die auf diese Theorie bezüglichen Aufzeichnungen der Scheda Ac sind in dem Abschnitt [V.] enthalten. Ehe wir auf diesen eingehen, bemerken wir noch wiederholt, daß die im art. [6.] S. 189 enthaltene Tabelle eine Erweiterung der am Schlusse des Abschnitts [III.], S. 183 befindlichen darstellt. Ferner möge noch mit wenigen Worten auf die Bedeutung der in den artt. [8.] und [9.] enthaltenen Aufzeichnungen hingewiesen werden, die auch schon im Bande III der Werke S. 423—425 abgedruckt sind\textsuperscript{**)}.
\par\indent In [31], S. 189 setzt \textsc{Gauss} für $M(1,\sin\varphi)$ eine trigonometrische Reihe an. Er scheint sich aber überzeugt zu haben, daß es auch hier zweckmäßiger ist, den reziproken Wert zu entwickeln, und findet, daß die Koeffizienten dieser Entwicklung in einfacher Weise von der lemniskatischen Periode abhängen. Insbesondere erweist sich das absolute Glied $N$ gleich $2\frac{\varpi^2}{\pi^2}$, ein Resultat, das mit dem im art. [13.] des Abschnitts [I.] der weiter unten abgedruckten Abhandlung »\textit{Zur Theorie der transcendenten Functionen gehörig}«\textsuperscript{***)} ausgesprochenen Theorem (das \textsc{Gauss} dort als »zu den Hauptsätzen dieser Theorie« gehörig bezeichnet) im Zusammenhang steht. In der Formel für den Maximalwert der dort betrachteten Funktion $\varphi a$ bedeutet nämlich $\varpi$ die lemniskatische Periode. Die im art. [9.] bewiesene Reihenidentität [41], S. 191
\Needspace{4\baselineskip}
\par\smallskip\noindent\rule{.25\linewidth}{.4pt}\par
{\small
\par\indent *) Das \textsc{Gauss}sche $z$ ist nichts anderes als das \textsc{Jacobi}sche $\sqrt q$.
\par\indent **) Der Abdruck ist im ganzen dem Original entsprechend, obgleich nicht philologisch genau. In der ersten Formel findet sich jedoch ein störendes Versehen, indem nämlich \textsc{Schering} statt des auf der rechten Seite von [40] stehenden Ausdrucks $P\cos\varphi$, gesetzt hat $\frac1{M(1,\cos\varphi)}$, während es nach [37] heißen müßte $\frac1{M(1,\sin\varphi)}$.
\par\indent ***) Aus dem Handbuch \textbf{16}, Bb; die Abhandlung ist auszugsweise bereits Werke III, S. 436 abgedruckt. Die hier in Rede stehende Stelle findet sich (mit Änderungen gegen die Handschrift) auf S. 441.
\par}
\clearpage
\fancyhead{}\fancyhead[R]{\small273}\fancyhead[C]{\scriptsize BEMERKUNGEN ZUM AGM. ABSCHNITT IV.}
\SourceLeaf{280}
\noindent hat \textsc{Gauss} im Jahre 1830 der Philosophischen Fakultät zu Göttingen als Preisaufgabe vorgeschlagen, die Aufgabe ist jedoch nicht gestellt worden\textsuperscript{*)}.
\par\indent Die Vergleichung der Abschnitte [II.], [III.], [IV.] zeigt vor allem, daß die Leisteaufzeichnungen [II.] früher entstanden sind als [III.] und [IV.]. Erstens sind die in [II.] angewandten Methoden sehr viel primitiver als die der Abschnitte [III.] und [IV.], nämlich rein formal, andererseits bilden die in [II.] gegebenen Reihenentwicklungen für den reziproken Wert des agM. und die daselbst zusammengestellten Formeln für den Ellipsenquadranten die Grundlage für die Betrachtungen von [III.] und [IV.]. Aber auch die Aufzeichnungen [II.] dürften nach dem 30. Mai 1799 entstanden sein; denn sie enthalten alle Hilfsmittel, um die in der Tagebuchnotiz Nr. 98 von diesem Tage (siehe S. 260) angegebene Gleichung
\[\frac\pi\varpi=M(\surd2,1)\]
\noindent zu beweisen, und es unterliegt auch keinem Zweifel, daß \textsc{Gauss} das auf dem Leisteschutzblatt bei dem Werte $1{,}198\pm$ später hinzugesetzte »$=\frac\pi\varpi$« (siehe oben S. 146, art. [3.] und die zugehörige Bemerkung S. 150) gerade um den 30. Mai 1799 geschrieben hat. Wahrscheinlich werden wir die Aufzeichnungen [II.] auf die Sommermonate 1799 anzusetzen haben.
\par\indent Für die Datierung der Aufzeichnungen [III.] und [IV.] bietet das \textit{Tagebuch} die nötigen Anhaltspunkte. Die Nr. 100 vom November 1799 »circa terminos medios arithmetico-geometricos multa nova deteximus« fällt zeitlich mit dem Beginn der Scheda Ac (November 1799, siehe S. 184 die Überschrift) zusammen; der erste Satz der am 14. Dezember 1798 zu Helmstedt niedergeschriebenen Nr. 101: »Medium arithmetico-geometricum tamquam quotientem duarum functionum transcendentium repraesentabile esse iam pridem inveneramus« bezieht sich offenbar auf die Gleichung [7], S. 186, und damit sind die artt. [1.]—[3.], oben S. 184—187, auf November 1799 datiert; sie sind noch in Braunschweig geschrieben. Der zweite Satz von Nr. 101 »nunc alteram harum functionum ad quantitates integrales reducibilem esse deteximus« bezieht sich auf die Gleichungen [13]—[16] des art. [4.], S. 187; diese Gleichungen und damit der ganze art. [4.] sind also schon im Dezember in Helmstedt geschrieben\textsuperscript{**)}. Die Funktion $\varphi(x)$ des Abschnitts [III.], S. 181, stimmt, wie schon bemerkt, mit $Q$ in Gl. [15], S. 187 überein; man liest ferner in der Nr. 102 des \textit{Tagebuchs} vom 23. Dezember: »Medium arithmetico-geometricum ipsum est quantitas integralis«, was mit der Bemerkung der Abhandlung von 1800, art. 8, Werke III, S. 370: »Hoc itaque modo media nostra arithmetico-geometrica ad quantitates integrales revocata sunt« fast wörtlich übereinstimmt. An dieser letzteren Stelle ist von der linearen Differentialgleichung zweiter Ordnung die Rede, der der reziproke Wert des agM. Genüge leistet; diese Differentialgleichung wird im Abschnitt [III.] aufgestellt, man wird also den in [III.] abgedruckten Zettel mit der Tagebuchnotiz Nr. 102 in Verbindung zu setzen haben.
\Needspace{5\baselineskip}
\par\smallskip\noindent\rule{.25\linewidth}{.4pt}\par
{\small
\par\indent *) Nach den Fakultätsakten lautet die von \textsc{Gauss} vorgeschlagene Aufgabe wie folgt:
\par\noindent »Demonstretur per methodos rigorosas atque concinnas aequalitas, quae intercedit inter seriem infinitam
\[1+\left(\frac12\right)^3x+\left(\frac{1\,.\,3}{2\,.\,4}\right)^3xx+\left(\frac{1\,.\,3\quad5}{2\,.\,4\,.\,6}\right)^3x^3+\left(\frac{1\,.\,3\,.\,5\,.\,7}{2\,.\,4\,.\,6\,.\,8}\right)^3x^4+\text{etc.}\]
\noindent atque quadratum seriei infinitae
\[1+\left(\frac14\right)^2x+\left(\frac{1\,.\,5}{4\,.\,8}\right)^2xx+\left(\frac{1\,.\,5\,.\,9}{4\,.\,8\,.\,12}\right)^2x^3+\left(\frac{1\,.\,5\,.\,9\,.\,13}{4\,.\,8\,.\,12\,.\,16}\right)^2x^4+\text{etc.«}\]
\noindent Wie man sieht, hat \textsc{Gauss} an die Stelle des in der Scheda Ac auftretenden $x^2$ hier $x$ gesetzt. Mit derselben Änderung ist die Gleichung [41] auch Werke III, S. 424 abgedruckt.
\par\indent **) Bei Gl. [13] beginnt ersichtlich andere Tinte.
\par}
\par\smallskip\noindent X1.\hfill35
\clearpage
\ArithmeticHead{274}{L}{ANALYSIS. NACHLASS.}
\SourceLeaf{281a}
\par\indent An den Schluß von [III.], S. 183, knüpft art. [5.] von [IV.], S. 188 unmittelbar an, ferner ist die Tabelle der Werte von $M\sin\varphi$ auf S. 189 eine wesentliche Erweiterung der am Ende von [III.], S. 183 befindlichen; \textsc{Gauss} wird also nach dem 23. Dezember 1799 wieder die Eintragungen in die Scheda Ac fortgesetzt und die artt. [5.]ff.\textsuperscript{*)} vor dem 13. Februar 1800 geschrieben haben. Die Tagebuchnotiz Nr. 103 von diesem Tage bezieht sich nämlich auf die Theorie der ternären quadratischen Formen, und eine von diesen Formen handelnde Eintragung findet sich auf S. 22 der Scheda (gedruckt Werke II, S. 311). Die vom 6. Mai 1800 an gemachten Tagebuchaufzeichnungen Nr. 105ff. stehen wieder mit den späteren Aufzeichnungen der Scheda Ac, zu deren Besprechung wir jetzt übergehen, in Übereinstimmung.
\Needspace{4\baselineskip}
\par\smallskip\noindent\rule{.25\linewidth}{.4pt}\par
{\small\par\indent *) Danach gehört also der Zettel [III.] sowohl der Zeit der Abfassung als auch dem Inhalte nach zwischen die artt. [4.] und [5.] des der Scheda Ac entnommenen Abschnitts [IV.]; daß er beim Abdruck vorangestellt worden ist, geschah nur, um die fortlaufende Wiedergabe des Inhalts der Scheda Ac nicht zu unterbrechen.\par}
\clearpage
\endgroup


